\begin{lemma}[Proved: integral Specht generators over a noetherian base]
\label{C:Specht_G0}
Let $k$ be any commutative noetherian ring and $n\geq0$.
The homomorphism
\[
 \bigoplus_{\lambda\vdash n}G_0(k)\longrightarrow G_0(k[\mathfrak S_n]),
 \qquad ([M_\lambda])_\lambda\longmapsto
       \sum_\lambda[S^\lambda_{\mathbb Z}\otimes_{\mathbb Z}M_\lambda]
\]
is surjective. Either the usual integral Specht lattices or their duals
may be used consistently.
\end{lemma}
\begin{proof}
We use the fieldwise input that the classes of Specht modules span
$G_0(F[\mathfrak S_n])$ for every field $F$, including characteristic two.
It follows from the unitriangular submatrix of the cellular decomposition
matrix, with diagonal entries one, as checked in the proof of
Proposition~\ref{C:DVR}.

Put $A=k[\mathfrak S_n]$; it is noetherian and finite free over its central
subring $k$. We prove the assertion by noetherian induction on the closed
support of a finite $A$-module $M$. Such a module admits a finite filtration
by $A$-submodules with each quotient annihilated by a prime of $k$: in a
nonzero quotient, choose a $k$-associated element $x$ whose annihilator is
a prime $\mathfrak p$ and take the $A$-submodule $Ax$. Centrality implies
that $\mathfrak p$ annihilates $Ax$, and noetherianity prevents an infinite
strictly increasing sequence of these submodules. Thus it suffices to treat
$M$ over the domain $k/\mathfrak p$ with support $V(\mathfrak p)$.

Let $F=\operatorname{Frac}(k/\mathfrak p)$. The fieldwise input writes
$[M_F]=\sum_\lambda a_\lambda[S^\lambda_F]$. This equality involves only
finitely many short exact sequences and isomorphisms of finite-dimensional
modules. Their matrices and relations can be spread out after inverting
one $f\notin\mathfrak p$; kernels and cokernels vanish after a further such
localization because they are finite. Therefore
\[
 [M]-\sum_\lambda a_\lambda[S^\lambda_{\mathbb Z}\otimes(k/\mathfrak p)]
\]
is a sum of classes of modules supported on
$V(\mathfrak p)\cap V(f)$, a proper closed subset of $V(\mathfrak p)$.
Here the elementary $G_0$ localization assertion can be verified directly:
finite modules over a central localization have finite lattices, two lattices
are commensurable up to a power of the inverted element, and exact sequences
are lifted by intersecting a lattice with the submodule and taking its
image in the quotient. Thus the kernel is generated by modules killed by
powers of $f$. The induction hypothesis treats these smaller supports.
The displayed Specht terms already belong to the asserted image, with
coefficient $[k/\mathfrak p]\in G_0(k)$. This completes the induction.
\end{proof}

\begin{theorem}[Proved: the complete answer for regular noetherian rings]
\label{C:regular_noetherian}
Let $k\ne0$ be a regular noetherian commutative ring. Then
$\vartheta_k$ is an isomorphism if and only if $k$ has characteristic zero
and the unit map $\mathbb Z\to K_0(k)$ is an isomorphism. No bound on the
Krull dimension is required. In particular it is an isomorphism for every
regular noetherian local ring of characteristic zero.
\end{theorem}
\begin{proof}
Necessity of the $K_0$ condition follows in degree zero, and positive
characteristic is excluded by Proposition~\ref{C:positive_characteristic}.
Assume the two conditions. Every finite $k$-module has a finite resolution
by finite projective $k$-modules. To recall why an unbounded global dimension
causes no gap here, fix a free resolution of that one module. At each prime,
regularity makes some syzygy free; the loci where a given finite syzygy is
locally free are open and increase with the syzygy index. They cover the
quasi-compact spectrum, so a single finite index works everywhere.

Apply the same construction using finite free $k[\mathfrak S_n]$-modules:
a sufficiently high syzygy is finite projective over $k$. Alternating sums
of these resolution terms give an inverse
$G_0(k[\mathfrak S_n])\to R_k(\mathfrak S_n)$ to the natural map.
Independence follows by iterated Schanuel comparison, and additivity by
horseshoe resolutions; the terminal syzygies are projective over $k$, so
their comparison sequences are admissible in $\mathcal R_k(\mathfrak S_n)$.
This is the same explicit Euler-characteristic argument as for the DVR,
with a finite number of additional syzygies.

By Lemma~\ref{C:Specht_G0}, the resulting group is generated by the classes
$S^\lambda_{\mathbb Z}\otimes M$, for $M$ finite over $k$. Resolving $M$
over $k$ expresses each such class by Specht lattices tensored with finite
projective $k$-modules. The hypothesis $K_0(k)=\mathbb Z$ then expresses
all of them as integer combinations of $[S^\lambda_k]$.
Burroughs's integral theorem expresses $[S^\lambda_{\mathbb Z}]$ as a
polynomial in the $\rho_i$. Its exact sequences are split over $\mathbb Z$,
so arbitrary base change to $k$ preserves them. Thus $\vartheta_k$ is
surjective in every degree.

Since $k$ has characteristic zero, $k\otimes\mathbb Q\ne0$ and has a
residue field $F$ of characteristic zero. Base change $k\to F$ gives a
map $R_k(\mathfrak S_*)\to R_F(\mathfrak S_*)$ carrying the $\rho_i$ to
the corresponding generators. The composite with $\vartheta_k$ is the
isomorphism of Proposition~\ref{C:Qalgebras} for the field $F$, so
$\vartheta_k$ is injective. This completes the proof.
\end{proof}

\begin{proposition}[Proved: passage to filtered unions and limits of rings]
\label{C:filtered_colimits}
For a filtered system of commutative rings $k_i$ with colimit $k$, and
any finite group $G$, the natural map
\[
 \mathop{\rm colim}_i R_{k_i}(G)\longrightarrow R_k(G)
\]
is an isomorphism. For symmetric groups these isomorphisms respect induction
products. Consequently a filtered colimit of rings satisfying C.1 also
satisfies C.1. This gives, for example, nonnoetherian mixed-characteristic
local examples
$\mathbb Z_p[t_1,t_2,\ldots]_{(p,t_1,t_2,\ldots)}$.
\end{proposition}
\begin{proof}
A finite projective module is represented by a finite idempotent matrix.
An action of a finite group, an equivariant homomorphism, or an equivariant
isomorphism is specified by finitely many additional matrices and polynomial
identities. Each therefore descends to some $k_i$, and two such data become
equal at some later stage if they become equal over $k$.
An admissible short exact sequence is split on the underlying projective
modules; its equivariance, exactness, and a chosen $k$-linear splitting are
again finitely many matrix identities and descend as well. The presentation
of $K_0$ by objects and these finite relations therefore commutes with the
filtered colimit. The explicit finite matrices for induction commute with
base change, proving the ring assertion. For the example, take the filtered
system $\mathbb Z_p[t_1,\ldots,t_r]_{(p,t_1,\ldots,t_r)}$, whose stages
satisfy Theorem~\ref{C:regular_noetherian}.
\end{proof}
